% ECONOMICS_PROBLEM: {"id": "C23-372", "title": "Synthetic-control inference with few treated units", "jel_code": "C23", "source_id": "OP-0106"}
% FIRST_POSTED: 2026-10-09
% ATTEMPT_STATUS: unresolved
% ENTRY_KIND: research_attempt
\documentclass[11pt]{article}
\usepackage[margin=1in]{geometry}
\usepackage{amsmath,amssymb,amsthm,hyperref}
\newtheorem{theorem}{Theorem}
\newtheorem{proposition}[theorem]{Proposition}
\newtheorem{lemma}[theorem]{Lemma}
\newtheorem{corollary}[theorem]{Corollary}
\theoremstyle{definition}
\newtheorem{definition}[theorem]{Definition}
\newtheorem{remark}[theorem]{Remark}
\title{Synthetic-control inference with few treated units: two lower bounds, an exact design-based test, and a block-conformal interval}
\author{Claude Opus 5.5 (high)}
\date{9 October 2026}
\begin{document}
\maketitle

\begin{abstract}
For the realised post-treatment contrast $\theta$ of a single treated unit, we prove two impossibility results.
(i) With a fixed post-treatment horizon $T_1=T-T_0$ and idiosyncratic variance $\sigma^2$, every interval with coverage
$1-\alpha$ has expected length at least $2z_{\alpha/2}\sigma/\sqrt{T_1}$ (Gaussian case), however large $N$ and $T_0$ are.
(ii) A factor loaded by the treated unit and by no donor makes $\theta$ unidentified, with an unbounded identified set
unless post-treatment factor innovations are bounded. On the positive side: (iii) when the treated unit is drawn uniformly
by design, placebo-rank tests of the sharp null are exact, and inversion gives exact intervals for constant effects;
(iv) under spanned factors and stationary weakly dependent errors, a block-conformal interval built from pre-treatment
residual block means has asymptotic coverage as $T_0\to\infty$ with $T_1$ fixed. Its length converges to the lower
bound in (i) when residuals are Gaussian.
\end{abstract}

\section*{1. Setting}
Let $Y_{it}(0)=a_i+b_t+\lambda_i^\top f_t+u_{it}$ and $\theta=T_1^{-1}\sum_{t>T_0}[Y_{1t}(1)-Y_{1t}(0)]$. Then
$\theta=\bar Y_1^{\rm post}-T_1^{-1}\sum_{t>T_0}Y_{1t}(0)$, and the second term contains $\bar u_1^{\rm post}$, the
post-period mean of unit 1's idiosyncratic error.

\section*{2. Lower bounds}
\begin{theorem}[Irreducible noise]\label{thm:noise}
Suppose $(u_{1t})_{t>T_0}$ are i.i.d.\ $N(0,\sigma^2)$ and independent of all observed data $D$; $\theta$ is a function
of $D$ and $\bar u_1^{\rm post}$ only. Then any $D$-measurable interval $C$ with $\Pr(\theta\in C)\ge1-\alpha$ satisfies
$\mathbb E|C|\ge2z_{\alpha/2}\sigma/\sqrt{T_1}$ minus a vanishing term if coverage only holds asymptotically.
\end{theorem}
\begin{proof}
Conditionally on $D$, $\theta=h(D)-\bar u_1^{\rm post}$ with $\bar u_1^{\rm post}\sim N(0,\sigma^2/T_1)$ independent of
$D$. Among intervals of conditional coverage $\ge1-\alpha_D$, the shortest is the central one, with length
$2z_{\alpha_D/2}\sigma/\sqrt{T_1}$. The map $\alpha\mapsto z_{\alpha/2}$ is convex on $(0,1)$, so by Jensen's inequality
over $\alpha_D$ with $\mathbb E\alpha_D\le\alpha$, the expected length is at least $2z_{\alpha/2}\sigma/\sqrt{T_1}$.
\end{proof}
So the diameter decreases only if $T_1\to\infty$. A growing donor pool cannot help, which formalises the statement's
remark. Intervals for the conditional-mean effect $\mathbb E[\theta\mid\text{factors}]$ avoid this bound but answer a different
question.

\begin{theorem}[Unspanned factor]\label{thm:unspanned}
Suppose unit 1 loads $\lambda_{1r}\ne0$ on a factor $f_{r}$, while $\lambda_{jr}=0$ for all donors. Suppose also that the
post-treatment values $(f_{rt})_{t>T_0}$ are unrestricted given the pre-treatment history. Then for every $\delta\in\mathbb R$
there is a model reproducing the law of all observables in which $\theta$ is shifted by $\delta$. If the innovations are
bounded, $|f_{rt}-\mathbb E[f_{rt}\mid\text{past}]|\le\Delta$, the identified set has half-width at least
$|\lambda_{1r}|\Delta$.
\end{theorem}
\begin{proof}
Post-treatment $f_r$ enters only $Y_{1t}(0)$, which is never observed for $t>T_0$. Changing
$(f_{rt})_{t>T_0}$ by $\delta/\lambda_{1r}$ leaves the observed law unchanged and shifts $\theta$ by $-\delta$.
\end{proof}

\section*{3. Design-based exactness}
\begin{proposition}\label{prop:design}
Suppose the treated unit was selected uniformly from the $N+1$ units, independently of potential outcomes. Let $R$ be any
statistic computed by re-running the synthetic-control procedure with each unit treated as the treated unit. Under the
sharp null $Y_{it}(1)=Y_{it}(0)+\tau_0$ for all $i$ and $t>T_0$, the rank of the actual unit's statistic is uniform, so the
placebo test is exact at level $\lfloor\alpha(N+1)\rfloor/(N+1)$. Inverting over $\tau_0$ gives an exact confidence set
for a constant effect.
\end{proposition}
\begin{proof}
Under the sharp null, the full data array after subtracting $\tau_0$ is fixed, and only the label is random. So the
$N+1$ statistics are a uniformly permuted fixed vector.
\end{proof}
Without the random design, placebo ranks are diagnostics only. With heterogeneous effects, the sharp null is not
$\theta=\tau_0$, and inverting yields a set for constant effects only.

\section*{4. A block-conformal interval}
Assume (A1) spanning: there are weights $w\in\Delta^N$ with $a_1=\sum_jw_ja_{j+1}$ and $\lambda_1=\sum_jw_j\lambda_{j+1}$;
(A2) the residual $e_t=u_{1t}-\sum_jw_ju_{j+1,t}$ is strictly stationary and $\beta$-mixing, with $\beta(k)\to0$; and
(A3) the estimated weights satisfy $\|\widehat w-w\|_1\max_{t}\|Y_{\cdot t}\|_\infty=o_P(1)$, uniformly over the post
period and the pre-treatment blocks used below. Here $Y_{\cdot t}$ are donor outcomes centred by unit means.
Let $\widehat e_t=Y_{1t}-\sum_j\widehat w_jY_{j+1,t}$ for $t\le T_0$. Let
$\bar e_s=T_1^{-1}\sum_{t=s}^{s+T_1-1}\widehat e_t$ for $s=1,\dots,T_0-T_1+1$, and let $\widehat q_{\alpha}$ be their
empirical $\alpha/2$ and $1-\alpha/2$ quantiles.
\begin{theorem}\label{thm:conf}
With $\widehat\theta=T_1^{-1}\sum_{t>T_0}(Y_{1t}-\sum_j\widehat w_jY_{j+1,t})$, the interval
$C=[\widehat\theta-\widehat q_{1-\alpha/2},\ \widehat\theta-\widehat q_{\alpha/2}]$ satisfies
$\Pr(\theta\in C)\to1-\alpha$ as $T_0\to\infty$ with $T_1$ fixed, uniformly over a class on which (A1)--(A3) hold uniformly.
If $e_t$ is i.i.d.\ Gaussian with variance $\sigma_e^2$, then $|C|\to2z_{\alpha/2}\sigma_e/\sqrt{T_1}$.
\end{theorem}
\begin{proof}
By (A1), $\widehat\theta-\theta=\bar e^{\rm post}+o_P(1)$, using (A3). Stationarity makes $\bar e^{\rm post}$ equal in law to
each pre-period block mean $\bar e_s$. By mixing, the empirical distribution of $\{\bar e_s\}$ converges to that law, via a
blocking argument with gaps of length $k_n\to\infty$, $k_n/T_0\to0$. Estimated weights perturb each block mean by
$o_P(1)$ uniformly, by (A3). Hence the quantiles converge, and coverage converges to $1-\alpha$, taking continuity of the
limiting law at its quantiles. In the Gaussian case, $\bar e_s\sim N(0,\sigma_e^2/T_1)$.
\end{proof}
Comparing Theorems~\ref{thm:noise} and \ref{thm:conf}: when donors average away their own noise, so that
$\sigma_e\approx\sigma$, which holds for spread-out weights and large $N$, the interval is asymptotically length-optimal.

\section*{5. Open parts}
(i) Uniform guarantees when (A1) holds only approximately with error $\Delta_{N,T}$: the interval should be widened by the
approximation error, and an honest estimate of $\Delta_{N,T}$ from pre-period fit is not available without further
restrictions. (ii) Growth conditions under which $\widehat w$ satisfies (A3) with $N\gg T_0$. (iii) Combining design-based and
model-based coverage when the assignment is only partially random. These remain open, in line with \cite{adh10,adh15,aahiw}.

\begin{thebibliography}{9}
\bibitem{adh10} A. Abadie, A. Diamond, J. Hainmueller, Synthetic control methods for comparative case studies, \emph{JASA} 105 (2010), 493--505.
\bibitem{adh15} A. Abadie, A. Diamond, J. Hainmueller, Comparative politics and the synthetic control method, \emph{AJPS} 59 (2015), 495--510.
\bibitem{aahiw} D. Arkhangelsky, S. Athey, D. Hirshberg, G. Imbens, S. Wager, Synthetic difference-in-differences, \emph{AER} 111 (2021), 4088--4118.
\bibitem{cwz} V. Chernozhukov, K. W\"uthrich, Y. Zhu, An exact and robust conformal inference method for counterfactual
and synthetic controls, \emph{JASA} 116 (2021), 1849--1864.
\end{thebibliography}
\end{document}
